\chapter{Glossary of Terms}
\label{app:glossary}

Plain-English definitions for security-specific terms used in this
thesis. Standard machine-learning terminology (precision, recall, ROC,
ETC.) is in the List of Abbreviations.

\begin{longtable}{p{4.0cm}p{11.0cm}}
    \toprule
    \textbf{Term} & \textbf{Definition} \\
    \midrule
    \endfirsthead
    \toprule
    \textbf{Term} & \textbf{Definition} \\
    \midrule
    \endhead
    Adversarial example & An input crafted to fool a machine-learning model into the wrong prediction. \\
    Adversarial training & Training a model on a mix of clean and adversarial examples to make it more robust. \\
    Anomaly detection & Detecting traffic that deviates from a learned baseline of normal behaviour. \\
    Bootstrap & A resampling-based statistical technique for estimating uncertainty. \\
    Brute force & Repeated guessing of credentials or keys until access is gained. \\
    Calibration & How closely a model's predicted probabilities match real-world frequencies. \\
    Concept drift & When the statistical properties of input data change over time. \\
    Day-split & Train on early days, test on later days; respects temporal order. \\
    DDoS & Distributed Denial of Service --- many sources flooding a target. \\
    EDR & Endpoint Detection and Response --- security software on every host. \\
    Evasion attack & An adversarial attack that flips a true positive to a false negative. \\
    Flow & A sequence of packets sharing source/destination IP, ports, and protocol. \\
    Heartbleed & A 2014 OpenSSL vulnerability that leaked memory contents. \\
    Kill chain & Lockheed Martin's seven-stage model of a typical intrusion. \\
    Leakage & When information from the test set inadvertently informs training. \\
    LODO & Leave-One-Day-Out --- train on $n - 1$ days, test on the held-out day. \\
    $L^\infty$ norm & The largest absolute change across all dimensions; bounds adversarial budgets. \\
    McNemar's test & A statistical test for comparing two classifiers on the same test set. \\
    MITRE ATT\&CK & An open framework cataloguing real-world attacker techniques. \\
    OODA loop & Boyd's Observe-Orient-Decide-Act decision cycle. \\
    Operating point & A specific threshold choice on the precision-recall trade-off curve. \\
    SHAP & A method for attributing a model's prediction to its input features. \\
    SIEM & Security Information and Event Management --- central alert aggregator. \\
    SOAR & Security Orchestration, Automation, and Response --- the playbook engine. \\
    SOC & Security Operations Centre --- the team that watches the network. \\
    Stratified split & A random split that preserves class proportions in every partition. \\
    Threat model & A formal description of who the attacker is and what they can do. \\
    Transferability & An adversarial example crafted against one model fooling another. \\
    Youden's J & TPR $-$ FPR; a prevalence-invariant threshold-selection criterion. \\
    \bottomrule
\end{longtable}
